\section{A noise-uniform readout transition}\label{sec:noisyreadout73}
The dependence on the readout basis can change scale when noise approaches
an extremal measurement. We determine that transition in a fixed
one-dimensional family, without a positive spectral margin uniform in
the noise. The proof is a finite-distance argument; no conclusion is
inferred merely from a Fisher-information calculation.

Let $\sigma_1,\sigma_2$ be the first two Pauli matrices. For
$\theta\in\R/(2\pi\Z)$ and a public visibility $0\le\lambda\le1$, put
\begin{equation}\label{eq:noisyfamily73}
 E_{\theta,\lambda}^{\pm}
 =\frac12\bigl(I\pm\lambda(\cos\theta\,\sigma_1+
                                  \sin\theta\,\sigma_2)\bigr),\qquad
 \mathcal M_{\theta,\lambda}^{\pm}(X)=\tr(E_{\theta,\lambda}^{\pm}X).
\end{equation}
Both outcome labels are retained. The output quantum system is
one-dimensional. In particular, $\theta$ and $\theta+\pi$ are not
identified by an outcome relabelling. Let $h(\theta,\phi)\in[0,\pi]$
be circular distance and set
\begin{equation}\label{eq:noisyscale73}
 \eta=1-\lambda^2,\qquad
 K_{N,\lambda}=\lambda\sqrt{N\min\{N,\eta^{-1}\}},\qquad N\ge1,
\end{equation}
where $\eta^{-1}=+\infty$ at $\lambda=1$. Thus $K_{N,1}=N$ and
$K_{N,0}=0$. The visibility may depend arbitrarily on $N$ and the
requested error; it is not a target coordinate charged in this theorem.

Write $\mathcal C_{N,\lambda}(\delta)$ for the least number of legal
memoryless binary instruments of this interface whose $d_N$-balls of
radius $\delta$ cover all targets in \eqref{eq:noisyfamily73}.
Centres need not belong to the displayed family. Define
$\mathcal C^{\mathrm{rat}}_{N,\lambda}(\delta)$ by restricting centres
to family members with rational $\cos\theta,\sin\theta$; rational
visibility then gives rational effects. A constant in the following
statement is independent of all three parameters $N,\lambda,\delta$.

\begin{theorem}[Noise-uniform finite-use coding]\label{thm:noisycoding73}
For all $N\ge1$, $0\le\lambda\le1$, and all pairs of angles,
\begin{equation}\label{eq:noisymetric73}
 \frac1{64}\min\{1,K_{N,\lambda}h(\theta,\phi)\}
 \le d_N(\mathcal M_{\theta,\lambda},\mathcal M_{\phi,\lambda})
 \le\min\{2,K_{N,\lambda}h(\theta,\phi)\}.
\end{equation}
There are absolute $c,C>0$ such that, for $0<\delta\le2^{-10}$,
\begin{equation}\label{eq:noisycover73}
 c\left(1+\frac{K_{N,\lambda}}\delta\right)
 \le\mathcal C_{N,\lambda}(\delta)
 \le\mathcal C^{\mathrm{rat}}_{N,\lambda}(\delta)
 \le C\left(1+\frac{K_{N,\lambda}}\delta\right).
\end{equation}
Consequently the optimum fixed-length reusable payload is
\begin{equation}\label{eq:noisybits73}
 \log_2\left(1+\frac{K_{N,\lambda}}\delta\right)+O(1).
\end{equation}
The upper construction is an exact rational algorithm on rational
visibility, tolerance and unit-circle input coordinates. At $\lambda=0$
the family is a singleton and its optimum payload is zero.
\end{theorem}
The constants and the additive remainder are uniform up to the erased
and projective endpoints. In contrast, they do not assert an optimal
leading constant, a large-error coherent covering law, or a coding law
in which an unknown visibility is also encoded.

\subsection{A two-point programme with its noise dependence retained}
\begin{lemma}[An adaptive upper modulus]\label{lem:noisyupper73}
Let $h=h(\theta,\phi)$ and $a=\lambda\sin(h/2)$. If $\eta>0$, then
\begin{equation}\label{eq:noisyfiniteupper73}
 d_N(\mathcal M_{\theta,\lambda},\mathcal M_{\phi,\lambda})
 \le\min\left\{2,\ 2Na,\
             2\sqrt{1-\left(\frac{\eta}{\eta+a^2}\right)^N}\right\}.
\end{equation}
For $\eta=0$ the bound $\min\{2,2Na\}$ remains valid.
\end{lemma}
\begin{proof}
The difference of the two positive-outcome effects has operator norm
$a$. On a bipartite input state, the two output blocks of the difference
are $B$ and $-B$, where
$B=\tr_A[((E^+_{\theta,\lambda}-E^+_{\phi,\lambda})\otimes I)\rho]$.
Duality against Hermitian contractions gives $\|B\|_1\le a$.
An eigenstate of the effect difference attains equality. Hence
$d_1=2a$, including reference assistance. Replacing the calls one at a
time and using trace-norm contraction proves $d_N\le2Na$ for every
common adaptive tester.

Choose angle representatives with difference $h$ and let $u,v$ be the
unit vectors along their midpoint and its oriented perpendicular in
the equatorial plane. The two Bloch vectors are $m\pm av$, where
$m=\lambda\cos(h/2)u$. Put
$b=(1-\|m\|^2)^{1/2}=(\eta+a^2)^{1/2}$.
The two Bloch vectors $m\pm bv$ have norm one, so they define two legal
binary projective instruments, say $\mathcal A_+$ and $\mathcal A_-$.
The target pair is exactly the two mixtures of these same instruments
with programme laws
\[
 p_\pm=\bigl((1\pm t)/2,(1\mp t)/2\bigr),\qquad t=a/b.
\]
This construction depends on the pair, which is permissible in a
metric upper bound. It is not one fixed processor for the whole circle.

For a fixed tester, first supply $N$ independent programme bits and,
at each requested call, apply the indicated $\mathcal A_+$ or
$\mathcal A_-$. All subsequent operations, including the reference,
feedback and a stopping rule, are a common channel of these bits.
Unused bits are discarded. The root fidelity of $p_+,p_-$ is
$\sqrt{1-t^2}$. Purifying the classical laws and using contraction gives
\[
 \|p_+^{\otimes N}-p_-^{\otimes N}\|_1
 \le2\sqrt{1-(1-t^2)^N}.
\]
This also follows directly from the classical fidelity inequality.
Substitution proves \eqref{eq:noisyfiniteupper73}.
Finally $1-(1-t^2)^N\le Nt^2$ and $a\le\lambda h/2$ imply
\[
 d_N\le\min\{N\lambda h,\sqrt{N/\eta}\,\lambda h,2\}.
\]
The limiting projective bound is supplied by the hybrid estimate,
not by division by zero.
\end{proof}
Classical simulation by common extremal channels and the effect of
noise on the Heisenberg scale are established methods
\cite{DKG67}. Environment/programme-state processing under adaptive
use is also an established framework \cite{DW72,WBHK72,WW72}.
The point needed here is the exact two-point expression
\eqref{eq:noisyfiniteupper73} and its uniform conversion to a family
cover, rather than a new claim about the origin of those methods.

\subsection{Finite entangled blocks and a binary product estimate}
\begin{lemma}[Symmetric Bernoulli products]\label{lem:bernoullimajority73}
Let $P_\pm$ be Bernoulli laws with success probabilities
$(1\pm a)/2$, where $0\le a\le1$. For every integer $m\ge1$,
\begin{equation}\label{eq:bernoullimajority73}
 \|P_+^{\otimes m}-P_-^{\otimes m}\|_1
                   \ge\tfrac12\min\{1,a\sqrt m\}.
\end{equation}
\end{lemma}
\begin{proof}
Retain the largest odd number $r=2j+1\le m$ of coordinates, so
$r\ge m/2$. For their majority event write
$g_r(a)=P_+(\text{majority})-P_-(\text{majority})$.
Differentiating the finite binomial sum and cancelling consecutive
terms gives
\[
 g_r'(a)=r\binom{2j}{j}4^{-j}(1-a^2)^j.
\]
For $j\ge1$, the bound $\binom{2j}{j}4^{-j}\ge1/(2\sqrt j)$
follows by induction: the ratio of consecutive left sides is
$(2j+1)/(2j+2)\ge\sqrt{j/(j+1)}$. For
$0\le a\le r^{-1/2}$, Bernoulli's inequality gives
$(1-a^2)^j\ge1-j/r\ge1/2$. Thus
$g_r'(a)\ge\sqrt r/(2\sqrt2)$ on that interval. The case $r=1$
is immediate. Integration from zero and monotonicity yield
$g_r(a)\ge(2\sqrt2)^{-1}\min\{1,a\sqrt r\}$.
Unhalved distance is at least twice this event difference; using
$r\ge m/2$ proves the stated bound.
\end{proof}

\begin{lemma}[A noise-uniform pairwise witness]\label{lem:noisylower73}
The lower inequality in \eqref{eq:noisymetric73} is witnessed by
nonadaptive entangled blocks, each of at most $N$ qubits, followed by
classical parity and majority processing.
\end{lemma}
\begin{proof}
A block of $k$ inputs is prepared in
$(|0\rangle^{\otimes k}+e^{i\alpha}|1\rangle^{\otimes k})/\sqrt2$.
Multiplying the $k$ output signs of \eqref{eq:noisyfamily73} has mean
$\lambda^k\cos(k\theta-\alpha)$. This identity follows by applying
$\bigotimes_{i=1}^k(E^+-E^-)$ to the two off-diagonal terms of the
block density matrix. For the two hypotheses choose the one common
phase $\alpha=k(\theta+\phi)/2+\pi/2$. Their parity laws then have
opposite biases of magnitude
\begin{equation}\label{eq:ghzbias73}
                        a_k=\lambda^k|\sin(kh/2)|.
\end{equation}
Repeating $m=\lfloor N/k\rfloor$ independent blocks and ignoring
remaining calls is an admissible tester. Its unhalved output distance
is at least $\tfrac12\min\{1,\sqrt m\,a_k\}$ by
Lemma~\ref{lem:bernoullimajority73}. The phase and block size may depend
on the pair, but never on the unknown hypothesis or a private outcome.

Assume $\lambda>0$ and put
$T=\min\{N,\eta^{-1}\}$, $k_0=\lfloor T\rfloor$. Then
$T/2\le k_0\le T$. For every $1\le k\le k_0$,
\begin{equation}\label{eq:visibilityblock73}
                              \lambda^k\ge\lambda/2.
\end{equation}
Indeed, when $0<\eta<1$, the inequality
$\log(1-\eta)\ge-\eta/(1-\eta)$ and
$k-1\le\eta^{-1}-1$ give $(k-1)\log\lambda\ge-1/2$.
At $\eta=0$ the assertion is immediate.

If $h\le1/k_0$, choose $k=k_0$. Since
$\sin(kh/2)\ge kh/\pi$ and $m\ge N/(2k)$, we obtain
\[
 \sqrt m\,a_k\ge
 \frac{\lambda\sqrt{Nk_0}}{2\pi\sqrt2}h
 \ge\frac{K_{N,\lambda}}{4\pi}h.
\]
This gives a lower constant $1/(8\pi)$.
If $k_0=1$, use $k=1$ for all $h\le\pi$ instead. Now $T<2$
unless $T=1$, and $\sqrt N\lambda\sin(h/2)\ge
K_{N,\lambda}h/(\pi\sqrt2)$; this is stronger than required.

It remains that $k_0\ge2$ and $h>1/k_0$. Here
$\lambda\ge1/\sqrt2$. Choose $k=\max\{1,\lfloor1/h\rfloor\}$.
Then $k\le k_0$ and $1/2\le kh\le\pi$. Equations
\eqref{eq:ghzbias73} and \eqref{eq:visibilityblock73} give
$a_k\ge1/(4\pi\sqrt2)$. Even one such block has distance at least
$1/(8\pi\sqrt2)>1/64$. The three cases prove the claim.
For $\lambda=0$ both sides of the desired inequality vanish.
\end{proof}

\subsection{Covering and exact rational realization}
\begin{proof}[Proof of Theorem~\ref{thm:noisycoding73}]
The metric assertion follows from the preceding lemmas. For a lower
cover choose points on the circle with circular separation at least
$256\delta/K_{N,\lambda}$ whenever that quantity is at most $\pi$.
An equally spaced packing has at least a fixed constant times
$K_{N,\lambda}/\delta$ points. Every two points have $d_N$-distance
at least $4\delta$, because $256\delta\le1/4$.
No radius-$\delta$ ball, even about an arbitrary legal instrument, can
contain two such points. If the required angular separation exceeds
$\pi$, the trivial lower bound one suffices, with a smaller absolute
constant. This proves the first inequality in \eqref{eq:noisycover73}.

For the upper bound use the two rational semicircle charts
\begin{equation}\label{eq:circlechart73}
 z_s(t)=s\left(\frac{1-t^2}{1+t^2},\frac{2t}{1+t^2}\right),
                 \qquad s\in\{-1,1\},\quad -1\le t\le1.
\end{equation}
Each chart has angular speed $2/(1+t^2)\le2$. Round $t$ to the
nearest $j/B$, $-B\le j\le B$, with a fixed rule at ties. The
circular error is at most $1/B$. Thus any integer
$B\ge\max\{1,K_{N,\lambda}/\delta\}$ yields a cover with at most
$4B+2$ centres. For rational inputs, an entirely rational choice is
\begin{equation}\label{eq:noisygrid73}
 B=\max\left\{1,\min\left\{
   \left\lceil\frac{4N\lambda}{\delta}\right\rceil,
   \left\lceil\sqrt{\frac{16N\lambda^2}{\eta\delta^2}}\right\rceil
                                  \right\}\right\},
\end{equation}
omitting the second entry when $\eta=0$. Integer squaring and rational
comparison compute the ceiling square root exactly. Both entries are
ceilings of the corresponding real bounds, so their minimum is
$\lceil4K_{N,\lambda}/\delta\rceil$.
For $z=(x,y)$ on the rational unit circle choose $s=1$ when $x\ge0$
and $s=-1$ otherwise, and set $t=sy/(1+sx)$. Its denominator is at
least one, so it is rational and belongs to $[-1,1]$.

The chart digit $(s,j)$ determines a legal effect pair with eigenvalues
$(1\pm\lambda)/2$. No approximate positivity test, irrational
normalization or search over real angles is needed. A fixed-length
index over $4B+2$ slots includes the chart sign. Its length gives
\eqref{eq:noisybits73}. Redundant chart endpoints cause at most a
constant overhead. At $\lambda=0$ dispense with the chart and emit the
single fair, input-erasing measurement.
\end{proof}

\begin{corollary}[A continuous family of horizon exponents]\label{cor:noisyexponents73}
Fix $0<\delta\le2^{-10}$. If $\lambda_N\to1$ and
$1-\lambda_N^2=N^{-\alpha+o(1)}$ for $\alpha\ge0$, then
\[
 \frac{\log_2\mathcal C_{N,\lambda_N}(\delta)}{\log_2N}
                       \longrightarrow\frac{1+\min\{\alpha,1\}}2.
\]
If $\lambda_N=N^{-\beta+o(1)}\to0$, with $\beta\ge0$, the limit is
$\max\{1/2-\beta,0\}$. The projective and erased endpoints have
coefficients one and zero, respectively.
\end{corollary}
\begin{proof}
Substitute the two regimes in \eqref{eq:noisyscale73} and use
\eqref{eq:noisycover73}. The term $1+K/\delta$ handles a vanishing
or subpolynomial scale without taking the logarithm of zero.
\end{proof}

\begin{corollary}[Noisy readout with conditional quantum states]\label{cor:noisypreparation73}
Fix $n$ and ranks $1\le r_+,r_-\le n$. Replace the scalar output in
\eqref{eq:noisyfamily73} by
\[
 \mathcal F_{\theta,\lambda,\rho}^{y}(X)
             =\tr(E_{\theta,\lambda}^{y}X)\rho_y,\qquad
 \tr\rho_y=1,\quad\operatorname{rank}\rho_y\le r_y,
\]
retaining the outcome $y$. Put
$V=\sum_{y\in\{+,-\}}[r_y(2n-r_y)-1]$.
For some fixed $\delta_0>0$, the optimum covering length is
\begin{equation}\label{eq:noisyprepbits73}
 \log_2\left(1+\frac{K_{N,\lambda}}\delta\right)
       +\frac V2\log_2N+V\log_2(1/\delta)+O(1),
 \qquad 0<\delta\le\delta_0.
\end{equation}
The constants are uniform in $N\ge1$ and $\lambda\in[0,1]$, at fixed
$n,r_+,r_-$. For rational visibility, upper centres can be rational and never
increase either conditional state rank; lower centres may be arbitrary legal
memoryless instruments of the same interface.
\end{corollary}
\begin{proof}
Discarding the quantum output leaves $\mathcal M_{\theta,\lambda}$,
so the readout lower bound is independent of the conditional states.
Conversely, using $I/2$ independently at each call produces exactly
$\Omega_\rho^{\otimes N}$, where
$\Omega_\rho=\tfrac12\rho_+\oplus\tfrac12\rho_-$, independently of
both $\theta$ and $\lambda$. Choose compact smooth patches of the two
state rank strata. Their product has dimension $V$ and is locally
bi-Lipschitz in the Hilbert--Schmidt metric on $\Omega_\rho$.
The local binary-test proof of Corollary~\ref{cor:localchoi68}, applied
to these states, supplies a packing of size
$c(\sqrt N/\delta)^V$ at output distance greater than $2\delta$ for
all sufficiently small fixed $\delta_0$. Cross it with the angular
packing used above. Two differing angles are separated by the readout
witness regardless of their states; pairs with the same angle are
separated by the maximally mixed-input witness. Triangle inequality
therefore proves the product lower bound for arbitrary centres.

For the upper cover first encode the angle to error at most $\delta/2$.
For fixed conditional states their preparation is a common operation
following the readout, so it cannot enlarge this error, including
within an adaptive tester. Encode each trace-one state separately by
Theorem~\ref{thm:factorcode68}, with preparation $N$-copy error at most
$\delta/4$. At each call a common processor can consume one copy of
each row state and select the state indexed by the measurement outcome.
Providing all $2N$ row programmes in advance controls every tester.
Triangle inequality on the two tensor products bounds their contribution
by $\delta/2$. The product of the two rank-aware factor codebooks has
size $C(\sqrt N/\delta)^V$. Rational state factors and the rational
angle charts give rational Choi blocks
$(E_{\theta,\lambda}^y)^{\mathsf T}\otimes\rho_y$ when $\lambda$
is rational. Each encoded state is trace one and rank nonincreasing,
which proves the asserted legality and payload bound.
\end{proof}

\begin{remark}[The retained flag and the scope of the transition]\label{rem:noisyscope73}
For $n=1$, forgetting $y$ yields the trace map for every $\theta$ and
$\lambda$; its covering number is one. For conditional states, forgetting
the flag likewise erases the readout whenever $\rho_+=\rho_-$.
Thus the observable interface is essential. At $0<\lambda<1$ the
measurement Choi blocks have full rank, whereas they have rank one at
$\lambda=1$. Their approach to that boundary is controlled explicitly
by $N(1-\lambda^2)$; a fixed-interior theorem alone would leave its
constants undetermined. The transition is not a general classification
of all coupled Choi tangent directions. Nor does a pairwise GHZ witness
supply a single globally successful estimator or extend this coherent
covering theorem to all fixed radii below two. Those are different
statements from \eqref{eq:noisycover73}.
\end{remark}
